\DocumentMetadata{
  pdfversion=2.0,
  pdfstandard=ua-2,
  lang=en-US,
  tagging=on,
  tagging-setup={math/setup=mathml-SE,tabsorder=structure}
}
\documentclass[11pt,letterpaper]{article}
\usepackage[margin=0.68in,includefoot]{geometry}
\usepackage{newcomputermodern}
\usepackage{amsmath,amscd,mathtools}
\usepackage{tikz,adjustbox}
\providecommand{\square}{\mdlgwhtsquare}
\usepackage[hidelinks]{hyperref}
\hypersetup{
  pdftitle={Topology PhD Exam, May 2019},
  pdfauthor={University of Florida Department of Mathematics},
  pdfsubject={Graduate mathematics examination},
  pdfdisplaydoctitle=true,
  bookmarksopen=true
}
\setcounter{secnumdepth}{0}
\setlength{\parindent}{0pt}
\setlength{\parskip}{0.28em}
\setlength{\emergencystretch}{3em}
\setlength{\leftmargini}{2.15em}
\setlength{\leftmarginii}{2.65em}
\setlength{\leftmarginiii}{3em}
\renewcommand{\labelenumi}{\textbf{\theenumi.}}
\pagestyle{empty}
\raggedbottom
\allowdisplaybreaks
\newenvironment{examproblems}[1]{%
  \begin{enumerate}%
  \setcounter{enumi}{#1}%
  \setlength{\topsep}{0.35em}%
  \setlength{\partopsep}{0pt}%
  \setlength{\itemsep}{0.48em}%
  \setlength{\parsep}{0.18em}%
}{\end{enumerate}}
\newenvironment{examsubparts}{%
  \begin{enumerate}%
  \setlength{\topsep}{0.18em}%
  \setlength{\partopsep}{0pt}%
  \setlength{\itemsep}{0.16em}%
  \setlength{\parsep}{0.1em}%
}{\end{enumerate}}
\newenvironment{examinstructions}{%
  \par\begingroup\itshape
}{%
  \par\endgroup\vspace{0.18em}
}
\makeatletter
\renewcommand\section{\@startsection{section}{1}{0pt}%
  {-1.25ex plus -.25ex minus -.1ex}%
  {0.55ex plus .1ex}%
  {\normalfont\large\bfseries}}
\renewcommand{\@maketitle}{%
  \newpage\null\vskip -1em%
  {\footnotesize\bfseries
    \href{https://www.ufl.edu/}{University of Florida}, \href{https://math.ufl.edu/}{Department of Mathematics}\par}%
  \vskip -0.5em%
  \tagstructbegin{tag=Title}%
    {\LARGE\bfseries\href{https://gma.math.ufl.edu/exams/topology-phd/}{Topology PhD Exam}, May 2019\par}%
  \tagstructend%
  \vskip 0.25em%
  \tagmcbegin{artifact}%
    \hrule height 1pt%
  \tagmcend%
  \vskip 1em%
}
\makeatother
\title{Topology PhD Exam, May 2019}
\author{}
\date{}
\begin{document}
\maketitle
\thispagestyle{empty}
\begin{examinstructions}
Work the following problems and show all your work. Support all statements to the best of your ability.
\end{examinstructions}
\begin{examproblems}{0}
\item
Prove that the topologist’s sine curve \(T=\{(x,\sin \frac{1}{x})\mid x\in(0,1]\}\cup\{(0,y)\mid y\in[0,1]\}\) is not path connected.
\item
Let~\(X\) be a connected completely regular topological space having more than one point. Can~\(X\) be countable?
\item
Show that the 2-sphere \(S^2\) is not a retract of the real projective plane \(\mathbb{RP}^2\), as well as that \(\mathbb{RP}^2\) is not a retract of \(S^2\).
\item
Consider the Klein Bottle, \(K\). Give its cellular chain complex and use it to calculate the homology groups of~\(K\) with integer coefficients.
\item
Let \(n>m\). Show that there is no map \(f:\mathbb{CP}^n\to\mathbb{CP}^m\) inducing a nontrivial map \(f^*:H^2(\mathbb{CP}^m;\mathbb{Z})\to H^2(\mathbb{CP}^n;\mathbb{Z})\).
\end{examproblems}
\begin{examinstructions}
Answer the following with complete definitions or statements or short proofs.
\end{examinstructions}
\begin{examproblems}{5}
\item
Draw a picture of the universal cover of the 2-sphere with a segment joining the north and south poles.
\item
Compute the Euler characteristic \(\chi(\mathbb{CP}^2\times\mathbb{RP}^2\times S^4)\).
\item
Give a definition of the compact-open topology for the space of functions.
\item
Construct a map \(f:T^2\to S^2\) of degree 3 where \(T^2=S^1\times S^1\) is a torus.
\item
Does every function \(f:\mathbb{N}\to S^1\) admit a continuous extension \(\bar f:\beta\mathbb{N}\to S^1\) to the Stone–Čech compactification?
\item
State the exact homology sequence for a pair of spaces.
\item
What can you say about the~\(k\)-th cohomology group of a closed orientable manifold of dimension~\(n\) for
\begin{examsubparts}
\item[\textnormal{(a)}]
\(k=n\)?
\item[\textnormal{(b)}]
\(k=n-1\)?
\end{examsubparts}
\item
State the Universal Coefficient Theorem for Cohomology.
\item
Show that two closed orientable surfaces of different genus are not homeomorphic.
\item
State the Tietze Extension Theorem.
\end{examproblems}
\end{document}
